\section{An actual gap before the bound $D<2$}
\label{v6gap:section}

The opposite-sign high-frequency event requires an actual gap around
zero, not the particular sufficient bound $D<2$. The stronger criterion
below uses the moment cell. It does not assert that any numerical cell
has acquired its hypotheses.

Let $X$ have mean zero, variance one and finite third absolute moment
$\rho$. Write $\gamma=\E X^3$, $m=\E X_+=\E X_->0$,
$\sigma=\sqrt{10}$ and
\[
 \tau(t)=\E[X\mathbf1_{\{X\ge t\}}],\quad
 J=\sum_{(u,v)}\tau_{u,v}(v-u)^2,\quad
 D=(3+\sigma)\rho-\gamma-3J,
\]
where the sum runs over bounded complementary intervals of the closed
support and $\tau_{u,v}$ is the constant value of $\tau$ on $(u,v)$.

\begin{lemma}[Moment-sensitive actual-gap event]\label{v6gap:event}
Suppose
\[
 1\le l\le\rho\le u,\quad D\le d,\quad d\ge0,\quad
 1<\sigma_-<\sqrt{10},\quad d<(\sigma_--1)l.
\]
Define the positive numbers
\[
 P_0=(\sigma_--1)l-d,\qquad S=l-\frac{d}{\sigma_--1}.
\]
Then the actual closed support has a central gap $(-A,B)$, with
$A,B>0$, and
\begin{equation}\label{v6gap:distance}
 (A+B)^2\ge 2S+\frac23P_0>2S>0.
\end{equation}
For independent copies of this same law, the actual event $E=\{XX'<0\}$
satisfies
\begin{equation}\label{v6gap:high-input}
 \Pp(E)\ge\alpha:=\frac1{2u^2},\qquad
 |X-X'|^2\ge2S+\frac23P_0\quad\hbox{on }E,\qquad
 \E[(X-X')^2\mid E]\le\frac2\alpha=4u^2.
\end{equation}
In particular no premise $d<2$ is needed for this event.
\end{lemma}
\begin{proof}
Signwise Tonelli gives $2\int|t|\tau(t)\,dt=\rho$.
For a gap on one side of zero, $(v-u)^2\le2\int_u^v|t|\,dt$.
If zero belonged to the closed support, all its gaps would be of that
kind, so $J\le\rho$. As $\gamma\le\rho$, this would imply
$D\ge(\sigma-1)\rho\ge(\sigma_--1)l>d$, a contradiction.
Closed support and both sign classes therefore give finite positive
inner endpoints $A,B$. Endpoint atoms are allowed.

Put
\[
 T_+=\E X_+^3-mB^2,\quad T_-=\E X_-^3-mA^2,\quad T=T_++T_-.
\]
These are nonnegative. The central gap contributes $m(A+B)^2$ to J;
write the remainder as $J_o$. For a positive outer gap,
$(v-u)^2\le2\int_u^v(t-B)\,dt$. Thus
\[
 J_o^+\le\E[X_+(X_+-B)^2]
       =T_+-2B(\E X_+^2-mB)\le T_+.
\]
The reflected argument gives $J_o^-\le T_-$, hence $J_o\le T$.
Expanding D exactly gives
\[
 D=m\{\sqrt{\sigma+1}A-\sqrt{\sigma-1}B\}^2
       +(\sigma-1)T_++(\sigma+1)T_-+3(T-J_o).
\]
Every term is nonnegative, so
$T\le D/(\sigma-1)\le d/(\sigma_--1)$ and
\[
 m(A^2+B^2)=\rho-T\ge S.
\]
For an additional product bound, summing the simpler gap inequality
gives $J\le\rho+2mAB$: the central gap exceeds
$2\int_{-A}^{B}|t|\,dt$ by exactly $2AB$.
Consequently
\[
 6mAB\ge\sigma\rho-\gamma-D
            \ge(\sigma-1)\rho-D\ge P_0.
\]
Since $m\le1/2$, combining the last two inequalities yields
\[
 (A+B)^2\ge\frac{S+P_0/3}{m}\ge2S+2P_0/3.
\]

There is no atom at zero. With $p=\Pp(X>0)$ and $q=1-p$,
the actual event has probability $2pq$. Conditional Cauchy gives
$1=\E X^2\ge m^2/p+m^2/q$, hence $pq\ge m^2$.
Also $1\le\E|X|\,\E|X|^3=2m\rho$, while
$2m=\E|X|\le1$. Therefore $2pq\ge2m^2\ge1/(2\rho^2)\ge\alpha$.
Opposite signs force $|X-X'|\ge A+B$. Finally
$\E(X-X')^2=2$, so conditioning on E gives at most
$2/\Pp(E)\le2/\alpha$.
\end{proof}

To use the complete high-frequency lemma, choose a positive rational
$d_0$ with $d_0^2\le2S+2P_0/3$, use event probability lower bound
$\alpha=1/(2u^2)$ and conditional second-moment ceiling $4u^2$, and
retain a valid fourth-moment ceiling for the same law. All the existing
window, damping and index gates remain. This licenses no automatic use
of the additional endpoint and moment conclusions bundled into the
older $d<2$ generic-event lemma.

In a finite ordered update, the new gap parameters must be obtained from
the preceding proved defect ceiling. Every complete high-frequency bound
then bounds the same actual Fourier integral. Its minimum with any
previous complete bound is valid. If the same low-frequency triple is
retained, the corresponding minima of the constant Fourier budgets are
valid as well. Raw and conserved-resource bounds must still be kept
distinct. This implication concerns the actual law throughout; it does
not replace the law by an artificial two-point kernel.
